\section{Method}
\label{sec:method}

\subsection{Problem formulation}

Let $X=\{x_t\}_{t=1}^{T}$ denote a multivariate industrial time series, where $x_t \in \mathbb{R}^{d}$ contains the sensor measurements at time step $t$. For each time index $t$, the model receives a past window
\begin{equation}
W_t = \{x_{t-L+1}, \ldots, x_t\}.
\end{equation}
The label belongs to the five-state set $\mathcal{S}=\{1,5,7,9,3\}$, corresponding to stop, normal operation, first-level degradation, second-level degradation, and fault occurrence.

The main task is fixed-horizon future-state classification:
\begin{equation}
\hat{y}_{t+\Delta} = f(W_t).
\end{equation}
This paper focuses on $\Delta=1$, which is the horizon supported by the main positive evidence. Longer-horizon results are reported only as a scope limitation.

\subsection{Overview}

SGTONetV6 is designed around a simple principle: common-state prediction and rare-state triggering should not be forced into the same flat classifier. The model therefore has two coupled paths. The first path is a conservative multiclass classifier for the future state. The second path is a rare-trigger branch that can override the base classifier only when transition constraints are satisfied.

\subsection{Patch temporal encoder}

The input window is divided into temporal patches and encoded into a shared representation. This gives the model local temporal context while keeping the encoder compact. The same encoder supports both the base future-state classifier and the rare-trigger branch.

\subsection{Conservative future-state classifier}

The base classifier predicts the five future states with a standard multiclass head. It is optimized to remain conservative under severe imbalance. This is important because an aggressive rare predictor can create many false alarms and reduce trust in deployment.

\subsection{Patch-attentive rare context}

The rare branch computes a patch-attentive context representation for the rare class. This mechanism is used because rare degradation evidence may appear only in a short part of the input window. Replacing this module with mean context substantially reduces rare-state recovery, as shown in the ablation study.

\subsection{Boundary-constrained rare trigger}

The rare head outputs a scalar rare score $r_t$. At inference time, the base prediction can be overridden only when all rare-trigger conditions are satisfied:
\begin{equation}
r_t \geq \tau,\quad b_t=1,\quad y_t \in \{5,7\},
\end{equation}
where $\tau$ is a calibrated threshold and $b_t$ is the transition-boundary flag. If the conditions are met, the prediction is changed to class 9. Otherwise, the model keeps the conservative base prediction. This rule is intentionally restrictive: it improves rare-state recall without allowing the rare head to fire globally.

\subsection{Threshold calibration and fallback prior}

The rare threshold is calibrated on validation data when possible. Because class 9 is extremely sparse, some validation splits may contain too few rare samples for stable calibration. SGTONetV6 therefore uses a fallback prior threshold. The ablation study shows that removing this fallback weakens class-9 recovery.
